\chapter{项目与模块}
\label{ch:18}

单文件适合实验，项目把入口、库模块、资源和构建配置组织在一起。模块决定声明的归属与可见性，配置决定编译器如何寻找源码和生成应用。

\section{先确定入口和目录}

把应用入口放在 src/main.chl，把业务逻辑按职责拆成模块，把网页等运行期资源放在独立目录。工作目录会影响相对文件路径，不能把 import 的源码搜索目录当成运行期资源目录。

\section{模块声明与限定名}

模块声明为名称提供归属。import 引入所需模块，限定名用于明确选择声明，尤其在多个模块出现同名类型时。不要依靠某一轮导入顺序解决重名；写出限定名并让诊断暴露真正冲突。

\par
\medskip
\begin{lstlisting}[style=chali]
inline int64 square(int64 value)
{
    return value * value;
}

void main()
{
    println("{}", square(3));
}
\end{lstlisting}
\par
\medskip

\par
\medskip
\begin{lstlisting}[style=out]
9
\end{lstlisting}
\par
\medskip

\section{项目配置与命令}

chali.toml 用于保存入口和项目选项。现有工程先检查实际配置和 chali help 输出，不把旧版 interpret 配置解释为当前存在的解释执行模式。当前 run 采用 checked AOT，aot 导出原生产物。工具链路径、目标信息和链接资产见第 24 章。

\section{可复现的工程边界}

在 README 记录启动目录、外部库、数据文件和目标平台。在 CI 中分开检查类型、运行例程与打包产物；生成文件从生成源维护，不通过编辑产物修补下一次必然被覆盖的内容。

\section{习题}

\begin{enumerate}

\item 为学生系统设计 src 与 web 目录。

\item 源码搜索与资源读取路径相同吗？

\item 两个模块都有 Error 时如何选择？

\item 给项目写出启动所需条件。

\end{enumerate}

\subsection*{参考解答与验收要点}

\begin{enumerate}

\item 按入口、HTTP、业务、存取及页面分工。

\item 不是，分别按 import 配置与运行期路径处理。

\item 使用模块限定名。

\item 记录编译器、目标、外部库、配置与工作目录。

\end{enumerate}

下一章：\hyperref[ch:19]{第 19 章}。

\section{番外：两个李逵，一个名字}

改写自《水浒传》李逵遇李鬼一节。

李逵行至岭边，跳出一个黑大汉，抡着两柄板斧，喝道：“留下买路钱！俺是黑旋风李逵！”李逵听得大怒，喝问：“你既是李逵，俺又是谁？”

二人争得面红耳赤，驿站伙计却只认短名，差点把工钱给了李鬼。掌柜拦住道：“名字相同，未必来自同一模块。且把限定名写全。”李逵写梁山，李鬼写树林，事情顿时清楚了许多。

李鬼又哭诉家有九十岁老母，声称自己只是一个本地同名声明，愿意遮蔽一下，绝不影响外部接口。李逵一时心软，许他暂留。回头再查，果见那厮在屋中商议如何冒领整个包的导出。

李逵一斧劈下，骂道：“遮蔽是名称解析的规矩，不是冒领的许可证！”

掌柜从此导入必查路径，同名必写归属。伙计问：“未声明模块者如何？”掌柜道：“记入小写 \texttt{program}，莫给它擅加一个大写头衔。”

李逵听罢，仍嫌字多。直到月底领钱，见收据上清清楚楚只付给梁山那位，才夸这套限定名虽不如板斧爽快，倒确实少一场人命官司。

\subsection*{技术对照}

模块名、文件搜索路径和符号名称共同决定解析。唯一短名便于调用，同名冲突时使用限定名；本地声明的遮蔽也需要明确检查。
